\documentclass{ximera}

\input{../../../preamble.tex}


\definecolor{fvdnavy}{HTML}{071D43}
\definecolor{fvdcyan}{HTML}{20BFC6}
\definecolor{fvdcyanlight}{HTML}{EFFCFB}
\definecolor{fvdmagenta}{HTML}{F10D69}
\definecolor{fvdmagentalight}{HTML}{FFF1F7}
\definecolor{fvdtext}{HTML}{163044}





%=================================================
% Pedagogische omgevingen
% PDF: tcolorbox
% HTML: learning-box
%=================================================

\newenvironment{voorkennis}
{%
    \ifdefined\HCode
        \HCode{%
<section class="learning-box learning-box--prior">
<div class="learning-box__header">
<span class="learning-box__icon" aria-hidden="true"></span>
<span class="learning-box__title">Voorkennis</span>
</div>
<div class="learning-box__content">
        }%
        \begin{itemize}
    \else
       \begin{tcolorbox}[
    enhanced,
    breakable,
    colback=fvdcyanlight,
    colframe=fvdcyan,
    coltext=fvdtext,
    boxrule=0.5pt,
    leftrule=4pt,
    arc=4mm,
    outer arc=4mm,
    left=7mm,
    right=6mm,
    top=5mm,
    bottom=5mm,
    before skip=6mm,
    after skip=6mm,
    title=Voorkennis,
    coltitle=fvdcyan,
    fonttitle=\bfseries\Large,
    attach title to upper={\par\smallskip},
    boxed title style={
        empty,
        left=0mm,
        right=0mm,
        top=0mm,
        bottom=0mm
    },
    overlay={
        \draw[
            fvdcyan,
            line width=0.5pt
        ]
        ([xshift=42mm,yshift=-13mm]frame.north west)
        --
        ([xshift=-7mm,yshift=-13mm]frame.north east);
    }
]
        \begin{itemize}
    \fi
}
{%
    \end{itemize}
    \ifdefined\HCode
        \HCode{%
</div>
</section>
        }%
    \else
        \end{tcolorbox}
    \fi
}


\newenvironment{leerdoelen}
{%
    \ifdefined\HCode
        \HCode{%
<section class="learning-box learning-box--goals">
<div class="learning-box__header">
<span class="learning-box__icon" aria-hidden="true"></span>
<span class="learning-box__title">Leerdoelen</span>
</div>
<div class="learning-box__content">
        }%
        \begin{itemize}
    \else
\begin{tcolorbox}[
    enhanced,
    breakable,
    colback=fvdmagentalight,
    colframe=fvdmagenta,
    coltext=fvdtext,
    boxrule=0.5pt,
    leftrule=4pt,
    arc=4mm,
    outer arc=4mm,
    left=7mm,
    right=6mm,
    top=5mm,
    bottom=5mm,
    before skip=6mm,
    after skip=6mm,
    title=Leerdoelen,
    coltitle=fvdmagenta,
    fonttitle=\bfseries\Large,
    attach title to upper={\par\smallskip},
    boxed title style={
        empty,
        left=0mm,
        right=0mm,
        top=0mm,
        bottom=0mm
    },
    overlay={
        \draw[
            fvdmagenta,
            line width=0.5pt
        ]
        ([xshift=42mm,yshift=-13mm]frame.north west)
        --
        ([xshift=-7mm,yshift=-13mm]frame.north east);
    }
]
        \begin{itemize}
    \fi
}
{%
    \end{itemize}
    \ifdefined\HCode
        \HCode{%
</div>
</section>
        }%
    \else
        \end{tcolorbox}
    \fi
}


\addPrintStyle{../../..}

\title{Oefeningen — Hoe verandert een functie?}
\author{Ingmar Herreman}

\begin{document}

% FVD_CHAPTER_DATA_START
% Automatisch gegenereerd uit index.tex.
% Niet handmatig aanpassen.
\fvdchapterdata{5}{Het gedrag van functies beschrijven}{3}
% FVD_CHAPTER_DATA_END




\fvdchapterdata{5}{Het gedrag van functies beschrijven}{3}

\maketitle
\FVDbeginExercises

\begin{opmerking}
De oefeningen met \(\blacklozenge\) zijn essentieel.
Omdat het leerplandoel op niveau \emph{analyseren} staat,
moet je veranderingen niet alleen herkennen maar ook verschillende
voorstellingen verbinden en conclusies motiveren.
\end{opmerking}


% ============================================================
% BASIS
% ============================================================

\FVDsection{Basis — verandering herkennen}

\begin{oefening}[Gelijke stappen]{\oefB\ESS}

Gegeven zijn drie tabellen.

\[
\begin{array}{c|ccccc}
x&0&1&2&3&4\\ \hline
f(x)&2&5&8&11&14
\end{array}
\]

\[
\begin{array}{c|ccccc}
x&0&1&2&3&4\\ \hline
g(x)&0&1&4&9&16
\end{array}
\]

\[
\begin{array}{c|ccccc}
x&0&1&2&3&4\\ \hline
h(x)&0&5&9&12&14
\end{array}
\]

\begin{question}
Bereken voor elke functie de opeenvolgende veranderingen.
\begin{oplossing}
Voor \(f\): \(+3,+3,+3,+3\).

Voor \(g\): \(+1,+3,+5,+7\).

Voor \(h\): \(+5,+4,+3,+2\).
\end{oplossing}
\end{question}

\begin{question}
Classificeer de stijging.
\begin{oplossing}
\(f\): constante stijging.

\(g\): toenemende stijging.

\(h\): afnemende stijging.
\end{oplossing}
\end{question}

\end{oefening}


\begin{oefening}[Drie manieren van dalen]{\oefB\ESS}

Classificeer telkens de daling.

\begin{question}
\[
\begin{array}{c|ccccc}
x&0&1&2&3&4\\ \hline
f(x)&12&10&8&6&4
\end{array}
\]

\begin{oplossing}
De veranderingen zijn \(-2,-2,-2,-2\).
Dit is constante daling.
\end{oplossing}
\end{question}

\begin{question}
\[
\begin{array}{c|ccccc}
x&0&1&2&3&4\\ \hline
g(x)&20&19&16&11&4
\end{array}
\]

\begin{oplossing}
De veranderingen zijn \(-1,-3,-5,-7\).
De daling wordt sterker: toenemende daling.
\end{oplossing}
\end{question}

\begin{question}
\[
\begin{array}{c|ccccc}
x&0&1&2&3&4\\ \hline
h(x)&20&14&10&8&7
\end{array}
\]

\begin{oplossing}
De veranderingen zijn \(-6,-4,-2,-1\).
De functie daalt nog steeds, maar steeds minder sterk:
afnemende daling.
\end{oplossing}
\end{question}

\end{oefening}


\begin{oefening}[Taal correct gebruiken]{\oefB\ESS}

Beoordeel de uitspraken.

\begin{question}
``Bij afnemende stijging daalt de functie.''
\begin{oplossing}
Fout. De functie stijgt nog steeds; alleen de stijging wordt minder sterk.
\end{oplossing}
\end{question}

\begin{question}
``Bij toenemende daling worden de functiewaarden kleiner.''
\begin{oplossing}
Juist. Bovendien wordt de afname steeds sterker.
\end{oplossing}
\end{question}

\begin{question}
``Een rechte met positieve richtingscoëfficiënt vertoont constante stijging.''
\begin{oplossing}
Juist.
\end{oplossing}
\end{question}

\end{oefening}


\begin{oefening}[Grafieken herkennen]{\oefB\ESS}

Koppel op \(x>0\) elk voorschrift aan de juiste beschrijving:

\[
f(x)=2x+1,\qquad
g(x)=x^2,\qquad
h(x)=\sqrt{x}.
\]

Kies uit:
\begin{enumerate}
  \item constante stijging;
  \item toenemende stijging;
  \item afnemende stijging.
\end{enumerate}

\begin{oplossing}
\[
f:\text{ constante stijging},
\qquad
g:\text{ toenemende stijging},
\qquad
h:\text{ afnemende stijging}.
\]
\end{oplossing}

\end{oefening}


% ============================================================
% VERDIEPING
% ============================================================

\FVDsection{Verdieping — analyseren en verbinden}

\begin{oefening}[Waarom moeten de stappen gelijk zijn?]{\oefU\ESS}

Een leerling vergelijkt:

\[
f(2)-f(1)=3
\]

met

\[
f(8)-f(3)=7
\]

en besluit: ``De stijging neemt toe, want \(7>3\).''

Analyseer de redenering.

\begin{oplossing}
De conclusie is niet verantwoord.
De eerste verandering hoort bij een stap van \(1\) in \(x\),
de tweede bij een stap van \(5\).
Om de verandering rechtstreeks te vergelijken, moeten we gelijke
stappen in de invoer gebruiken of de verandering per eenheid vergelijken.
\end{oplossing}

\end{oefening}


\begin{oefening}[Van beschrijving naar schets]{\oefU\ESS}

Schets een mogelijke grafiek die:
\begin{itemize}
  \item voor \(x<0\) toenemend stijgt;
  \item voor \(0<x<3\) afnemend stijgt;
  \item voor \(x>3\) toenemend daalt.
\end{itemize}

\begin{question}
Wat moet er rond \(x=3\) met het gewone verloop gebeuren?
\begin{oplossing}
De functie gaat van stijgend naar dalend.
Daar ontstaat dus een lokaal maximum.
\end{oplossing}
\end{question}

\begin{question}
Is de grafiek uniek bepaald?
\begin{oplossing}
Nee. De kenmerken leggen het kwalitatieve gedrag vast,
maar niet alle functiewaarden of de exacte vorm.
\end{oplossing}
\end{question}

\end{oefening}


\begin{oefening}[Opwarming analyseren]{\oefU\ESS}

De temperatuur van een voorwerp wordt om de minuut gemeten:

\[
\begin{array}{c|rrrrrr}
t\text{ (min)}&0&1&2&3&4&5\\ \hline
T\text{ (}^{\circ}\mathrm C\text{)}&20&38&51&60&66&70
\end{array}
\]

\begin{question}
Toon met de tabel aan dat de temperatuur stijgt.
\begin{oplossing}
Elke volgende temperatuur is groter dan de vorige.
\end{oplossing}
\end{question}

\begin{question}
Classificeer de stijging.
\begin{oplossing}
De toenames zijn
\[
+18,\ +13,\ +9,\ +6,\ +4.
\]
Ze worden kleiner. Er is dus afnemende stijging.
\end{oplossing}
\end{question}

\begin{question}
Interpreteer dit in de context.
\begin{oplossing}
Het voorwerp wordt nog steeds warmer, maar de temperatuur neemt per minuut
steeds minder toe. De opwarming vertraagt.
\end{oplossing}
\end{question}

\end{oefening}


\begin{oefening}[Afstand en beweging]{\oefU\ESS}

Voor een voertuig wordt gemeten:

\[
\begin{array}{c|rrrrrr}
t\text{ (s)}&0&1&2&3&4&5\\ \hline
s\text{ (m)}&0&3&8&15&24&35
\end{array}
\]

\begin{question}
Bereken de afstandstoenames per seconde.
\begin{oplossing}
\[
3,\ 5,\ 7,\ 9,\ 11\text{ m}.
\]
\end{oplossing}
\end{question}

\begin{question}
Wat kun je besluiten over \(s(t)\)?
\begin{oplossing}
De afstand stijgt en de stijging neemt toe.
\end{oplossing}
\end{question}

\begin{question}
Wat betekent dit kwalitatief voor de beweging?
\begin{oplossing}
Per seconde wordt meer afstand afgelegd dan in de vorige seconde.
Het voertuig beweegt dus steeds sneller in de beschouwde meetgegevens.
\end{oplossing}
\end{question}

\end{oefening}


\begin{oefening}[Vind de begripsfout]{\oefU\ESS}

Een leerling bekijkt een grafiek die stijgt maar steeds vlakker wordt en zegt:

``De functie begint te dalen, want de stijging neemt af.''

Leg precies uit wat fout is.

\begin{oplossing}
De leerling verwart de functiewaarden met de verandering ervan.
De functiewaarden nemen nog steeds toe, dus de functie stijgt.
Alleen de hoeveelheid waarmee ze toenemen wordt kleiner.
Dat is afnemende stijging, niet daling.
\end{oplossing}

\end{oefening}


\begin{oefening}[Vier situaties vergelijken]{\oefU\ESS}

Geef voor elke beschrijving aan welke situatie bedoeld wordt.

\begin{enumerate}
  \item De functiewaarden nemen toe en de grafiek wordt steiler.
  \item De functiewaarden nemen af en de grafiek vlakt af.
  \item De functiewaarden nemen toe en de grafiek vlakt af.
  \item De functiewaarden nemen af en de grafiek wordt steiler.
\end{enumerate}

\begin{oplossing}
\begin{enumerate}
  \item toenemende stijging;
  \item afnemende daling;
  \item afnemende stijging;
  \item toenemende daling.
\end{enumerate}
\end{oplossing}

\end{oefening}


% ============================================================
% GEVORDERD
% ============================================================

\FVDsection{Gevorderd — onderzoeken en modelleren}

\begin{oefening}[Eén functie, verschillende regimes]{\oefG}

Een functie heeft achtereenvolgens:
\[
\text{toenemende stijging}
\longrightarrow
\text{afnemende stijging}
\longrightarrow
\text{afnemende daling}
\longrightarrow
\text{toenemende daling}.
\]

Maak een kwalitatieve schets en duid aan:
\begin{itemize}
  \item waar een lokaal extremum noodzakelijk ontstaat;
  \item waar het soort stijging/daling verandert zonder dat een extremum nodig is.
\end{itemize}

\begin{oplossing}
Tussen afnemende stijging en afnemende daling verandert de functie van stijgend
naar dalend; daar is een lokaal maximum noodzakelijk.

Tussen toenemende en afnemende stijging kan de functie aan beide kanten blijven stijgen.
Daar hoeft dus geen extremum te liggen. Hetzelfde geldt voor een verandering
tussen twee soorten daling.
\end{oplossing}

\end{oefening}


\begin{oefening}[Welke tabel past?]{\oefG}

Ontwerp zelf drie tabellen met \(x=0,1,2,3,4,5\):
\begin{enumerate}
  \item één met toenemende stijging;
  \item één met afnemende stijging;
  \item één met afnemende daling.
\end{enumerate}

Zorg ervoor dat de functiewaarden niet eenvoudigweg afkomstig zijn van de
standaardvoorbeelden \(x\), \(x^2\) of \(\sqrt{x}\).
Motiveer telkens met de opeenvolgende verschillen.

\begin{oplossing}
Meerdere antwoorden zijn mogelijk. Essentieel is dat de verschillen bij gelijke
stappen respectievelijk positieve en groter wordende, positieve en kleiner wordende,
en negatieve maar in absolute waarde kleiner wordende waarden hebben.
\end{oplossing}

\end{oefening}


\begin{oefening}[GeoGebra: functies vergelijken]{\oefG}

Onderzoek op \(x\geq0\):

\[
f(x)=x,\qquad
g(x)=x^2,\qquad
h(x)=\sqrt{x}.
\]

Maak in GeoGebra naast de grafieken ook een tabel voor
\(x=0,1,2,3,4,5\).

\begin{enumerate}
  \item Classificeer de stijging van elke functie.
  \item Onderbouw elke conclusie zowel grafisch als numeriek.
  \item Onderzoek vervolgens \(-f(x)\), \(-g(x)\) en \(-h(x)\).
  \item Formuleer wat de spiegeling in de \(x\)-as doet met
        stijgen/dalen en met toenemend/afnemend veranderen.
\end{enumerate}

\begin{oplossing}
\(f\) heeft constante stijging, \(g\) toenemende stijging en \(h\)
afnemende stijging op het beschouwde gebied.

Spiegeling in de \(x\)-as maakt de functiewaarden tegengesteld:
constante stijging wordt constante daling,
toenemende stijging wordt toenemende daling,
en afnemende stijging wordt afnemende daling.
\end{oplossing}

\end{oefening}


% ============================================================
% MINIMUM
% ============================================================


\end{document}
